\section{Introduction}
\label{sec:intro}

Chain-of-thought (CoT) reasoning has become the default strategy for improving LLM performance on complex tasks~\citep{wei2022cot}. Large reasoning models such as OpenAI's o1~\citep{openai2024o1} and DeepSeek-R1~\citep{deepseek2025r1} generate extended internal reasoning traces spanning thousands of tokens. The prevailing assumption is straightforward: more reasoning leads to better answers. But does this assumption hold when models encounter tasks outside their training distribution?

Recent work challenges the ``more is always better'' narrative. \citet{wu2025more} establish an inverted U-shaped curve between CoT length and accuracy on in-distribution (ID) tasks, showing that excessively long traces accumulate errors. \citet{su2025underthinking} demonstrate that incorrect responses are dramatically longer than correct ones, while \citet{yang2025longtoshort} reveal that training-free length reduction degrades behavioral consistency even when accuracy is preserved. These findings raise a natural question: \emph{how does the length--accuracy relationship change when models face out-of-distribution tasks, and can we exploit this relationship for more robust deployment?}

Despite growing interest in efficient reasoning~\citep{sui2025overthinking, yue2025efficient}, no prior work has systematically measured how the reasoning length--accuracy curve shifts across ID versus OOD benchmarks, or whether uncertainty signals can guide adaptive length selection for OOD robustness. \citet{aggarwal2025l1} show that RL-trained length control generalizes to some OOD tasks, but do not characterize the shape of the length--robustness curve. \citet{valmeekam2025reasonless} find that semantically corrupted traces can match or exceed genuine traces on OOD tasks, suggesting that trace content and length may be decoupled from actual reasoning quality.

We address this gap with a systematic investigation of reasoning trace length as a proxy for robustness. Using budget-forcing prompts, we control reasoning length at five levels (from zero reasoning to unconstrained) in GPT-4.1 and evaluate on four benchmarks spanning different domains and difficulty levels: MATH~\citep{hendrycks2021math} (ID, competition math), GSM8K~\citep{cobbe2021gsm8k} (easy OOD), MMLU-Pro~\citep{wang2023mmlupro} (knowledge-heavy OOD), and MuSR~\citep{sprague2023musr} (hard multi-step reasoning OOD). We measure accuracy, token usage, confidence calibration, and the \robgap---the ID--OOD accuracy difference---as a function of trace length.

Our results reveal a nuanced picture. On MATH and MMLU-Pro, accuracy increases monotonically with trace length (46\%$\rightarrow$78\% and 63\%$\rightarrow$79\%, respectively), with diminishing returns beyond ${\sim}$500 tokens. On GSM8K, accuracy is flat at ${\sim}$93\% regardless of budget. On MuSR, however, we observe the hypothesized non-monotonic pattern: short reasoning (5.1\%) is \emph{worse} than no reasoning (10.3\%), accuracy peaks at moderate length (21.8\%), and declines slightly at the longest budget (20.5\%). Most strikingly, model confidence is well-calibrated ID ($\rho{=}0.63$ on MATH) but completely uninformative OOD ($\rho{\approx}0$ on MuSR), causing a confidence-based adaptive strategy to fail precisely where it is most needed.

We make the following contributions:
\begin{itemize}[leftmargin=*,itemsep=0pt,topsep=0pt]
    \item We conduct the first systematic study of how the reasoning trace length--accuracy relationship varies across ID and OOD benchmarks, revealing that the relationship is task-dependent rather than universally non-monotonic.
    \item We identify the ``underthinking trap''---a novel finding that forcing brief reasoning can be \emph{worse} than no reasoning at all on hard OOD tasks, with short CoT underperforming no-CoT by 5.2 percentage points on MuSR.
    \item We demonstrate that model-expressed confidence is catastrophically miscalibrated on hard OOD tasks ($\rho{\approx}0$), undermining confidence-based adaptive length selection and revealing a fundamental limitation of self-supervised uncertainty estimation for reasoning control.
    \item We evaluate an adaptive two-pass strategy that achieves 32\% token savings on ID tasks while matching fixed-budget accuracy, but fails on OOD tasks due to overconfidence---highlighting the need for external calibration signals.
\end{itemize}
